\section{Test-Time Scaling：从 LLM 到扩散模型}

\subsection{预训练时代的终结与推理时计算}

Ilya Sutskever（OpenAI 联合创始人）曾提出一个影响深远的观点：预训练时代正在走向终结。对于大语言模型（LLM）而言，互联网上可用的训练数据已经接近用尽。即使继续扩大模型规模或增加训练计算量，性能提升也将越来越有限。

\begin{knowledgebox}{Test-Time Scaling 的兴起}
Test-time scaling 的核心思想是：与其在训练时投入更多计算资源，不如在推理时花费更多计算来提升输出质量。这一思想在 LLM 领域已经取得了显著成功：
\begin{itemize}
    \item OpenAI 的 O1/O3 模型通过``思考模式''在推理时花费更多时间来改善答案
    \item DeepSeek R1 通过 RL 训练模型产生中间思考 token，实现了``aha moment''——模型能自我发现并纠正推理错误
    \item 在 ARC-AGI 基准测试中，O3 的分数从此前模型的 10\%--20\% 跃升至 88\%
\end{itemize}
\end{knowledgebox}

\subsection{从 LLM 到扩散模型的迁移}

课程的核心问题是：能否将 test-time scaling 的思想从 LLM 迁移到扩散模型和流模型？讲者认为这是一个极具前景的方向，且相比训练时计算，推理时技术对计算资源的要求更低，这使得小规模研究团队也有机会开发出性能优异的模型。

\begin{importantbox}{Proposer-Verifier 框架}
Test-time scaling 的基本架构由两个组件构成：
\begin{itemize}
    \item \textbf{Proposer（提议者）}：预训练的生成模型（如扩散模型），负责生成候选输出。
    \item \textbf{Verifier（验证者）}：另一个模型或函数，评估 proposer 的输出质量。可以是神经网络、代码验证器或任何评分函数。
\end{itemize}
基于 verifier 的反馈，可以（1）引导 proposer 的生成过程，或（2）让 proposer 生成多个候选，从中选择最优。这与 GAN 的``生成器-判别器''框架异曲同工。
\end{importantbox}

\subsection{RL 视角下的推理时引导}

Dan（OpenAI 工程师）的演讲生动阐述了 test-time scaling 的愿景：就像 1907 年的爱因斯坦需要8年才能发展出广义相对论，而 O3 模型只需要一分钟的``思考''就能解答广义相对论期末考试题——这依赖于推理时花费更多计算。

Yann LeCun 曾将预训练比作``大蛋糕''，强化学习只是``蛋糕上的樱桃''。但 OpenAI 的路线图表明，未来 RL 计算将在总计算中占据越来越大的比重，最终可能远超预训练计算。

\begin{warningbox}{研究趋势的启示}
讲者观察到，当前 AI 研究社区已不再仅仅关注开发单一新技术或发表论文。越来越多的研究团队（包括创业公司）正在利用 test-time scaling 技术，用相对较少的计算资源（相比从头训练大模型）来超越现有模型的性能。这为资源有限的学术实验室和小团队提供了重要机会——不需要拥有数千张 GPU，也可以通过巧妙的推理时技术取得竞争力强的结果。
\end{warningbox}

\subsection{本章小结}

Test-time scaling 已经在 LLM 领域展现了巨大潜力。将这一思想迁移到扩散模型，通过 proposer-verifier 框架在推理时提升生成质量，是当前研究的重要方向。

